\subsection{Momentum placement and mobility}
\label{sec:monotone-mobility}

The direct update stores the scalar \(q^n\) but no previous increment. Momentum
carries previous increments into the current step. Euclidean heavy-ball momentum \cite{polyak1964} retains a fraction
\(\beta\) of the previous parameter increment \(s^n\),
\begin{equation}
  s^{n+1}=\beta s^n-\eta g^n,\qquad
  \theta^{n+1}=\theta^n+s^{n+1},\qquad 0\leq\beta<1,
  \label{eq:heavy-ball-step}
\end{equation}
and reduces to gradient descent at \(\beta=0\).

\subsubsection{Placement of momentum}

Consider first computing the direct PB--SAV increment and then appending
heavy-ball memory:
\begin{subequations}
\label{eq:appended-momentum}
\begin{align}
  \left(\eta^{-1}I_d+B_\alpha^n\right)\delta_{\mathrm{PB}}^{n+1}
  &=-\frac{q^n}{Q^n}g^n,
  \label{eq:appended-base-step}\\
  s^{n+1}&=\delta_{\mathrm{PB}}^{n+1}+\beta s^n,\\
  \theta^{n+1}&=\theta^n+s^{n+1},
  \qquad q^{n+1}=q^n+d^{n+1},
\end{align}
\end{subequations}
where \(d^{n+1}=(g^n)^{\mathsf T}s^{n+1}/(2Q^n)\). The scalar update uses
the total increment. Define the Euclidean modified energy
\begin{equation}
  \Hcal_{\mathrm{app}}^n=(q^n)^2+\frac{\norm{s^n}^2}{2\eta}.
  \label{eq:appended-energy}
\end{equation}

\begin{proposition}
\label{prop:appended-momentum-defect}
The update \cref{eq:appended-momentum} satisfies
\begin{equation}
\begin{aligned}
  \Hcal_{\mathrm{app}}^n-\Hcal_{\mathrm{app}}^{n+1}
  ={}&\frac{\norm{s^{n+1}-\beta s^n}^2}{2\eta}
  +\frac{1-\beta^2}{2\eta}\norm{s^n}^2\\
  &+(d^{n+1})^2+\alpha\norm{s^{n+1}}_{S_m^n}^2
  -\beta(s^{n+1})^{\mathsf T}B_\alpha^n s^n.
\end{aligned}
\label{eq:appended-momentum-defect}
\end{equation}
The last term has no fixed sign.
\end{proposition}

\begin{proof}
The total increment obeys
\((\eta^{-1}I_d+B_\alpha^n)(s^{n+1}-\beta s^n)=-(q^n/Q^n)g^n\).
Take its inner product with \(s^{n+1}\) and apply
\cref{eq:common-identities}. Subtracting the change in kinetic energy and
using Euclidean polarization gives \cref{eq:appended-momentum-defect}.
\end{proof}

\subsubsection{The momentum and the mobility}

We place momentum on the right-hand side of the implicit solve, so that the
pullback acts on the gradient and memory together. This removes the extra
\(\beta B_\alpha^n s^n\) term in the appended state equation.
Let \(p^n\in\R^d\) be the momentum and let the mobility
\(M^n\in\R^{d\times d}\), symmetric positive definite, map it to the
parameter increment \(M^np^n\); this is the mobility in gradient-flow
terminology and a variable-metric preconditioner in optimization. The new
momentum and increment are related by
\begin{equation}
  \Delta\theta^{n+1}=M^{n+1}p^{n+1},
  \qquad
  p^{n+1}=(M^{n+1})^{-1}\Delta\theta^{n+1}.
  \label{eq:covector-map}
\end{equation}
We assume that the mobility is nonincreasing in the Loewner order:
\begin{equation}
  0\prec M^{n+1}\preceq M^n.
  \label{eq:abstract-mobility-order}
\end{equation}
The PB--SAV momentum update is
\begin{equation}
  \left((M^{n+1})^{-1}+\eta B_\alpha^n\right)
  \Delta\theta^{n+1}
  =\beta p^n-\eta\frac{q^n}{Q^n}g^n,
  \label{eq:abstract-momentum-step}
\end{equation}
or, equivalently,
\begin{equation}
  \frac{p^{n+1}-\beta p^n}{\eta}
  +B_\alpha^n\Delta\theta^{n+1}
  +\frac{q^n}{Q^n}g^n=0.
  \label{eq:abstract-momentum-balance}
\end{equation}
With Euclidean mobility, omitting the pullback term and setting \(q^n/Q^n=1\)
in this balance gives the heavy-ball recursion \cref{eq:heavy-ball-step},
with \(p^n=s^n\).
The parameter and scalar states remain
\begin{equation}
  \theta^{n+1}=\theta^n+\Delta\theta^{n+1},
  \qquad
  q^{n+1}=q^n+d^{n+1}.
  \label{eq:abstract-momentum-states}
\end{equation}

\subsubsection{The inertial quadratic model}
\label{sec:momentum-filtering}

The implicit solve \cref{eq:abstract-momentum-step} has a quadratic
optimization interpretation.

\begin{proposition}
\label{prop:inertial-model}
The increment \(\Delta\theta^{n+1}\) solving
\cref{eq:abstract-momentum-step} is the unique minimizer of
\begin{equation}
  \mathcal Q(s)
  =\frac{q^n}{Q^n}(g^n)^{\mathsf T}s
  +\frac12 s^{\mathsf T}B_\alpha^ns
  +\frac{1}{2\eta}\norm{s-\beta M^{n+1}p^n}_{(M^{n+1})^{-1}}^2 .
  \label{eq:inertial-quadratic-model}
\end{equation}
Appending the momentum after the solve gives instead
\begin{equation}
  \Delta_{\mathrm{app}}^{n+1}
  =-\eta\frac{q^n}{Q^n}
    \bigl((M^{n+1})^{-1}+\eta B_\alpha^n\bigr)^{-1}g^n
   +\beta M^{n+1}p^n,
  \label{eq:frozen-appended-comparator}
\end{equation}
and the two increments satisfy
\begin{align}
  \Delta_{\mathrm{app}}^{n+1}-\Delta\theta^{n+1}
  &=\beta\eta\bigl((M^{n+1})^{-1}+\eta B_\alpha^n\bigr)^{-1}
    B_\alpha^nM^{n+1}p^n,
  \label{eq:appended-model-step-defect}\\
  \mathcal Q(\Delta_{\mathrm{app}}^{n+1})-\mathcal Q(\Delta\theta^{n+1})
  &=\frac{1}{2\eta}
  \norm{\Delta_{\mathrm{app}}^{n+1}-\Delta\theta^{n+1}}
       _{(M^{n+1})^{-1}+\eta B_\alpha^n}^2\geq0.
  \label{eq:appended-model-excess}
\end{align}
\end{proposition}

\begin{proof}
The Hessian of \(\mathcal Q\) is
\(B_\alpha^n+\eta^{-1}(M^{n+1})^{-1}\succ0\), and its optimality equation is
\cref{eq:abstract-momentum-step}. Subtracting the two increments and using
\(\bigl((M^{n+1})^{-1}+\eta B_\alpha^n\bigr)M^{n+1}
 =I_d+\eta B_\alpha^nM^{n+1}\) gives
\cref{eq:appended-model-step-defect}. Expanding \(\mathcal Q\) about its
minimizer gives \cref{eq:appended-model-excess}.
\end{proof}

\begin{remark}
\label{rem:inertial-gauss-newton}
For a fixed \(\alpha\), let \(\Psi_\alpha(\theta)\) stack
\(\sqrt{2(1-\alpha)}\,Q(\theta)\) on \(\sqrt{\alpha}\,\psi(\theta)\), so that
\(\tfrac12\norm{\Psi_\alpha}^2=F+C\), its Gauss--Newton matrix is
\(B_\alpha\), and its gradient is \(g\). At exact
tracking \(q^n=Q^n\), the shifted model \(\mathcal Q(s)+F+C\) is a
Gauss--Newton model for the residual map \(\Psi_\alpha\), together with a
proximal term centered at \(\beta M^{n+1}p^n\) and measured in
\(\norm{\cdot}_{(M^{n+1})^{-1}}\); otherwise the linear term is scaled by
\(q^n/Q^n\). Levenberg--Marquardt methods for neural networks also combine curvature with
adaptive momentum \cite{pooladzandi2022improving}, without the scalar
coupling that produces the energy law below.
\end{remark}

To see how the solve acts on memory, let \(b_j\geq0\) be the eigenvalues of
\((M^{n+1})^{1/2}B_\alpha^n(M^{n+1})^{1/2}\) and let \(u_j\) be the
coordinates of \((M^{n+1})^{1/2}p^n\) in a corresponding orthonormal
eigenbasis. In those coordinates the solve multiplies the \(j\)th component
of the stored momentum by \(\beta/(1+\eta b_j)\), and
\begin{equation}
  \mathcal Q(\Delta_{\mathrm{app}}^{n+1})-\mathcal Q(\Delta\theta^{n+1})
  =\frac{\beta^2\eta}{2}
  \sum_j\frac{b_j^2}{1+\eta b_j}u_j^2 .
  \label{eq:spectral-memory-coefficient}
\end{equation}

\begin{remark}
\label{rem:memory-damping}
Directions with larger \(b_j\) damp the stored momentum more strongly, and
where \(b_j=0\) the coefficient remains \(\beta\). The excess in
\cref{eq:spectral-memory-coefficient} is measured in the quadratic model, not
in the modified energy of \cref{thm:monotone-mobility-dissipation}.
\end{remark}


\subsubsection{Curvature action of the momentum step}
\label{sec:curvature-action}

The following comparison measures the approximation error of the correction
through its action on the reference increment.

\begin{proposition}
\label{prop:curvature-action}
Let \(M\succ0\), let \(A,B\succeq0\) be symmetric, and let \(\eta>0\).
For a common right-hand side \(f\in\R^d\), define
\begin{equation}
  (M^{-1}+\eta B)\Delta_B=f,
  \qquad
  (M^{-1}+\eta A)\Delta_A=f.
  \label{eq:curvature-reference-steps}
\end{equation}
Then
\begin{equation}
  \Delta_B-\Delta_A
  =\eta(M^{-1}+\eta B)^{-1}(A-B)\Delta_A,
  \label{eq:curvature-action-identity}
\end{equation}
and
\begin{equation}
  \norm{\Delta_B-\Delta_A}_{M^{-1}}
  \leq\eta\norm{M^{1/2}(A-B)\Delta_A}.
  \label{eq:curvature-action-bound}
\end{equation}
In particular, \(\Delta_B=\Delta_A\) if and only if
\((A-B)\Delta_A=0\).
\end{proposition}

\begin{proof}
Subtracting the two systems gives \cref{eq:curvature-action-identity}.
Since \(M^{1/2}BM^{1/2}\succeq0\), we have
\(\norm{(I_d+\eta M^{1/2}BM^{1/2})^{-1}}_2\leq1\); multiplying
\cref{eq:curvature-action-identity} by \(M^{-1/2}\) then gives
\cref{eq:curvature-action-bound}. The equality criterion follows from
invertibility of \(M^{-1}+\eta B\).
\end{proof}

\begin{remark}
\label{rem:curvature-action}
For the momentum update, take \(M=M^{n+1}\), \(B=B_\alpha^n\), and
\(f=\beta p^n-\eta(q^n/Q^n)g^n\); a Gauss--Newton matrix, or a positive
semidefinite Hessian at the current state, can serve as \(A\). A low-rank \(B\)
approximates \(\Delta_A\) when the weighted action error on the right-hand
side of \cref{eq:curvature-action-bound} is small, even when \(A-B\) is not
small in operator norm, and reproduces it exactly when
\((A-B)\Delta_A=0\). The quadratic and regression tests
in \cref{sec:quadratic-experiment,sec:regression-experiment} measure the
matrix error and the step error separately.
\end{remark}
